% ========== SECURITY & PERMISSIONS ==========
% Security architecture, permission model, trust verification, sandboxing & isolation
% Source: Symphony/Content/Security Framework, Trust Verification, Symphony Permissions Integration Guide

\section*{Security \& Permissions}
\addcontentsline{toc}{section}{Security \& Permissions}

\lettrine{S}{ecurity} in Symphony is not an afterthought but a foundational architectural principle that permeates every layer of the system. The security architecture implements defense-in-depth strategies, fine-grained permission models, and comprehensive trust verification to ensure safe operation in enterprise environments while maintaining the flexibility required for AI-driven development workflows.

\subsection*{Security Architecture}
\addcontentsline{toc}{subsection}{Security Architecture}

Symphony's security architecture is built on the principle of least privilege, with multiple layers of protection that work together to provide comprehensive security coverage.

\subsubsection*{Multi-Layer Security Model}

The security architecture implements protection at multiple levels to ensure comprehensive coverage against various threat vectors.

\textbf{Layer 1: Process Isolation}
\begin{compactlist}
    \item \textbf{Extension Sandboxing}: Each extension runs in isolated processes with restricted capabilities
    \item \textbf{Resource Limits}: CPU, memory, and I/O limits enforced at the OS level
    \item \textbf{Namespace Isolation}: Separate namespaces for filesystem, network, and IPC resources
    \item \textbf{Capability Dropping}: Extensions run with minimal required capabilities
\end{compactlist}

\textbf{Layer 2: Language-Level Security}
\begin{compactlist}
    \item \textbf{Memory Safety}: Rust's compile-time guarantees prevent buffer overflows and use-after-free
    \item \textbf{Type Safety}: Strong typing prevents type confusion attacks
    \item \textbf{Safe Concurrency}: Rust's ownership model prevents data races and concurrent access issues
    \item \textbf{Bounds Checking}: Automatic bounds checking prevents array overflow attacks
\end{compactlist}

\textbf{Layer 3: Application-Level Security}
\begin{compactlist}
    \item \textbf{Permission System}: Fine-grained permissions for resource access control
    \item \textbf{Input Validation}: Comprehensive input sanitization and validation
    \item \textbf{Cryptographic Protection}: End-to-end encryption for sensitive data
    \item \textbf{Audit Logging}: Comprehensive logging for security monitoring and forensics
\end{compactlist}

\textbf{Layer 4: Network Security}
\begin{compactlist}
    \item \textbf{TLS Encryption}: All network communication encrypted with TLS 1.3
    \item \textbf{Certificate Pinning}: Certificate validation to prevent man-in-the-middle attacks
    \item \textbf{Network Segmentation}: Isolated network segments for different security zones
    \item \textbf{Firewall Integration}: Integration with host and network firewalls
\end{compactlist}

\subsubsection*{Threat Model}

Symphony's security architecture addresses a comprehensive threat model covering various attack vectors and scenarios.

\begin{table}[h]
\centering
\caption{Threat Categories and Mitigations}
\label{tab:threat-model}
\begin{tabular}{@{}llll@{}}
\toprule
\textbf{Threat Category} & \textbf{Risk Level} & \textbf{Primary Mitigation} & \textbf{Secondary Mitigation} \\
\midrule
Malicious Extensions & High & Sandboxing & Permission System \\
Code Injection & High & Input Validation & Memory Safety \\
Data Exfiltration & Medium & Access Control & Audit Logging \\
Privilege Escalation & High & Least Privilege & Process Isolation \\
Network Attacks & Medium & TLS Encryption & Certificate Pinning \\
Supply Chain & Medium & Code Signing & Trust Verification \\
\bottomrule
\end{tabular}
\end{table}

\subsection*{Permission Model}
\addcontentsline{toc}{subsection}{Permission Model}

Symphony implements a sophisticated permission model that balances security with usability, providing fine-grained control over resource access while maintaining developer productivity.

\subsubsection*{Role-Based Access Control (RBAC)}

The foundation of Symphony's permission system is a flexible RBAC implementation that supports complex organizational structures.

\textbf{Core RBAC Components}:
\begin{expandedlist}
    \item \textbf{Users}: Individual developers with unique identities and authentication credentials
    \item \textbf{Roles}: Named collections of permissions that define specific access patterns
    \item \textbf{Permissions}: Atomic access rights to specific resources or operations
    \item \textbf{Resources}: Protected entities including files, extensions, AI models, and system functions
    \item \textbf{Contexts}: Environmental factors that influence permission evaluation
\end{expandedlist}

\textbf{Built-in Roles}:
\begin{compactlist}
    \item \textbf{Developer}: Standard development permissions with project access
    \item \textbf{Senior Developer}: Enhanced permissions including extension installation
    \item \textbf{Team Lead}: Team management and project configuration permissions
    \item \textbf{Administrator}: System administration and security configuration
    \item \textbf{Security Officer}: Security monitoring and audit access
\end{compactlist}

\subsubsection*{Capability-Based Security}

Beyond RBAC, Symphony implements capability-based security for fine-grained access control.

\textbf{Capability Types}:
\begin{compactlist}
    \item \textbf{File System Capabilities}: Read, write, execute permissions for specific paths
    \item \textbf{Network Capabilities}: Network access permissions for specific hosts and ports
    \item \textbf{Extension Capabilities}: Permissions to install, configure, and execute extensions
    \item \textbf{AI Model Capabilities}: Access to specific AI models and computational resources
    \item \textbf{System Capabilities}: Access to system functions and configuration
\end{compactlist}

\textbf{Capability Delegation}:
\begin{compactlist}
    \item \textbf{Temporary Delegation}: Time-limited capability grants for specific operations
    \item \textbf{Scoped Delegation}: Capabilities limited to specific resources or contexts
    \item \textbf{Revocable Delegation}: Capabilities that can be revoked at any time
    \item \textbf{Auditable Delegation}: All capability delegations are logged and auditable
\end{compactlist}

\subsubsection*{Dynamic Permission Evaluation}

Symphony's permission system evaluates access requests dynamically based on current context and system state.

\textbf{Evaluation Factors}:
\begin{compactlist}
    \item \textbf{User Identity}: Authenticated user identity and associated roles
    \item \textbf{Resource Sensitivity}: Classification level of the requested resource
    \item \textbf{Operation Type}: Type of operation being requested (read, write, execute)
    \item \textbf{Time Context}: Time-based restrictions and business hours policies
    \item \textbf{Location Context}: Network location and geographic restrictions
    \item \textbf{System State}: Current system load and security posture
\end{compactlist}

\begin{infobox}[title=Permission Evaluation Algorithm]
\textbf{Step 1}: Authenticate user identity and validate session

\textbf{Step 2}: Determine applicable roles and capabilities

\textbf{Step 3}: Evaluate resource-specific permissions

\textbf{Step 4}: Apply contextual restrictions and policies

\textbf{Step 5}: Log access decision and grant or deny access
\end{infobox}

\subsection*{Trust Verification}
\addcontentsline{toc}{subsection}{Trust Verification}

Trust verification ensures that all code and extensions running within Symphony are authentic, unmodified, and from trusted sources.

\subsubsection*{Code Signing Infrastructure}

Symphony implements a comprehensive code signing infrastructure to verify the authenticity and integrity of all executable code.

\textbf{Signing Process}:
\begin{expandedlist}
    \item \textbf{Developer Certificates}: Developers obtain certificates from trusted Certificate Authorities
    \item \textbf{Code Signing}: All extensions and updates are digitally signed before distribution
    \item \textbf{Timestamp Signing}: Signatures include trusted timestamps to prevent replay attacks
    \item \textbf{Certificate Validation}: Signatures are validated against current certificate revocation lists
    \item \textbf{Trust Chain Verification}: Complete certificate chain validation to root CA
\end{expandedlist}

\textbf{Signature Verification}:
\begin{compactlist}
    \item \textbf{Installation Time}: All extensions verified before installation
    \item \textbf{Runtime Verification}: Periodic verification of loaded extensions
    \item \textbf{Update Verification}: All updates verified before application
    \item \textbf{Integrity Checking}: Regular integrity checks of installed extensions
\end{compactlist}

\subsubsection*{Trust Levels}

Symphony implements a multi-level trust system that allows for graduated trust based on source and verification status.

\begin{table}[h]
\centering
\caption{Trust Levels and Permissions}
\label{tab:trust-levels}
\begin{tabular}{@{}llll@{}}
\toprule
\textbf{Trust Level} & \textbf{Source} & \textbf{Verification} & \textbf{Permissions} \\
\midrule
Verified Publisher & Official Store & Full Signature & All Capabilities \\
Trusted Developer & Known Publisher & Valid Signature & Standard Capabilities \\
Community Verified & Community Store & Community Review & Limited Capabilities \\
Self-Signed & Local Developer & Self-Signed Cert & Minimal Capabilities \\
Unsigned & Unknown Source & No Verification & Sandbox Only \\
\bottomrule
\end{tabular}
\end{table}

\subsubsection*{Reputation System}

Symphony maintains a reputation system that tracks the behavior and reliability of extensions and developers.

\textbf{Reputation Factors}:
\begin{compactlist}
    \item \textbf{Security History}: Track record of security issues and responses
    \item \textbf{Quality Metrics}: User ratings, bug reports, and performance metrics
    \item \textbf{Community Standing}: Community contributions and peer reviews
    \item \textbf{Compliance Record}: Adherence to security policies and guidelines
\end{compactlist}

\subsection*{Sandboxing \& Isolation}
\addcontentsline{toc}{subsection}{Sandboxing \& Isolation}

Symphony's sandboxing and isolation mechanisms provide strong security boundaries while maintaining the performance and functionality required for AI-driven development.

\subsubsection*{Extension Sandboxing}

Each extension runs in a carefully controlled sandbox environment that limits its access to system resources.

\textbf{Sandbox Implementation}:
\begin{compactlist}
    \item \textbf{Process Isolation}: Each extension runs in a separate process with restricted privileges
    \item \textbf{Filesystem Isolation}: Chroot-like isolation with access only to permitted directories
    \item \textbf{Network Isolation}: Network access restricted to explicitly permitted destinations
    \item \textbf{IPC Isolation}: Inter-process communication limited to approved channels
\end{compactlist}

\textbf{Resource Limits}:
\begin{compactlist}
    \item \textbf{CPU Limits}: CPU usage capped to prevent resource exhaustion
    \item \textbf{Memory Limits}: Memory allocation limits with automatic cleanup
    \item \textbf{I/O Limits}: Disk and network I/O rate limiting
    \item \textbf{Time Limits}: Execution time limits for long-running operations
\end{compactlist}

\subsubsection*{AI Model Isolation}

AI models require special isolation considerations due to their computational requirements and potential security implications.

\textbf{Model Sandbox Features}:
\begin{compactlist}
    \item \textbf{GPU Isolation}: Dedicated GPU memory allocation with cleanup guarantees
    \item \textbf{Model Encryption}: Models encrypted at rest and in memory when possible
    \item \textbf{Input Sanitization}: All model inputs sanitized to prevent injection attacks
    \item \textbf{Output Filtering}: Model outputs filtered to prevent data exfiltration
\end{compactlist}

\textbf{Inference Security}:
\begin{compactlist}
    \item \textbf{Prompt Injection Protection}: Detection and prevention of prompt injection attacks
    \item \textbf{Output Validation}: Validation of model outputs for safety and appropriateness
    \item \textbf{Resource Monitoring}: Continuous monitoring of model resource usage
    \item \textbf{Anomaly Detection}: Detection of unusual model behavior or outputs
\end{compactlist}

\subsubsection*{Data Isolation}

Sensitive data requires special protection to prevent unauthorized access or exfiltration.

\textbf{Data Classification}:
\begin{compactlist}
    \item \textbf{Public}: Non-sensitive data with no access restrictions
    \item \textbf{Internal}: Internal data with standard access controls
    \item \textbf{Confidential}: Sensitive data with enhanced protection
    \item \textbf{Restricted}: Highly sensitive data with maximum protection
\end{compactlist}

\textbf{Protection Mechanisms}:
\begin{compactlist}
    \item \textbf{Encryption at Rest}: All sensitive data encrypted when stored
    \item \textbf{Encryption in Transit}: All data transfers encrypted with TLS
    \item \textbf{Memory Protection}: Sensitive data cleared from memory after use
    \item \textbf{Access Logging}: All access to sensitive data logged and monitored
\end{compactlist}

\begin{successbox}
\textbf{Security Performance}: Symphony's security architecture is designed to provide strong protection with minimal performance impact. Most security operations add less than 5\% overhead to normal operations, while providing enterprise-grade security guarantees.
\end{successbox}

\subsubsection*{Compliance and Auditing}

Symphony's security architecture supports compliance with major security frameworks and regulations.

\textbf{Compliance Support}:
\begin{compactlist}
    \item \textbf{SOC 2}: System and Organization Controls compliance
    \item \textbf{ISO 27001}: Information Security Management System compliance
    \item \textbf{GDPR}: General Data Protection Regulation compliance
    \item \textbf{HIPAA}: Health Insurance Portability and Accountability Act compliance
\end{compactlist}

\textbf{Audit Capabilities}:
\begin{compactlist}
    \item \textbf{Comprehensive Logging}: All security-relevant events logged with timestamps
    \item \textbf{Tamper-Proof Logs}: Cryptographically signed logs to prevent tampering
    \item \textbf{Real-Time Monitoring}: Continuous monitoring with alerting for security events
    \item \textbf{Forensic Analysis}: Tools for detailed forensic analysis of security incidents
\end{compactlist}

\begin{alertbox}
\textbf{Security Considerations}: While Symphony's security architecture provides strong protection, security is an ongoing process that requires regular updates, monitoring, and response to emerging threats. Organizations should implement appropriate security policies and procedures to complement Symphony's technical security measures.
\end{alertbox}